\documentclass[10pt,a4paper]{amsart}
\usepackage[margin=19mm]{geometry}
\usepackage{amsmath,amssymb,mathtools}
\usepackage[scr=rsfs,cal=euler]{mathalfa}
\usepackage{xfrac}
\usepackage[unicode,hidelinks,bookmarks=false]{hyperref}
\newcommand{\ie}{i.e.\ }
\newcommand{\cf}{cf.\ }
\newcommand{\resp}{respectively\ }
\newcommand{\Subsection}{\subsection}
\providecommand{\nobreakdash}{\nobreak\hbox{-}}
\setlength{\emergencystretch}{2em}
\multlinegap0pt
\usepackage{bm}
\usepackage{bbm}
\usepackage{dsfont}
\usepackage{xspace}
\usepackage{cancel}
\usepackage{mathtools}
\usepackage{amsthm}
\makeatletter
\@ifundefined{theorem}{\newtheorem{theorem}{Theorem}}{}
\makeatother
\newcommand{\innpr}[2]{{\left\langle {#1},{#2}\right\rangle}} %inner product
\newcommand{\mb}[1]{\boldsymbol{#1}} % For Greek letters
\newcommand{\m}[1]{\mathbf{#1}} % General matrix notation
\newcommand{\rd}{\mathrm{d}} %Integral
\newcommand{\s}[1]{\mathscr{#1}} % General integral Operator Notation
\newcommand{\vertiii}[1]{\mathopen{|||} #1 \mathclose{|||}}
\renewcommand{\b}[1]{\mathbb{#1}} % General Field notation or RKHS
\renewcommand{\c}[1]{\mathcal{#1}} % Generic space notation
\title[Kernel Embedding for Operator-Valued Measures and Its Applic]{Kernel Embedding for Operator-Valued Measures and Its Application to Quantum Tomography}
\author{Philipp Nikolas Mayer, Ho Yun}
\date{Source: arXiv:2605.25146v1 (CC BY 4.0)}
\begin{document}
\maketitle

\noindent\emph{Source and licence.} Kernel Embedding for Operator-Valued Measures and Its Application to Quantum Tomography. Version arXiv:2605.25146v1 (CC BY 4.0), distributed under the Creative Commons Attribution 4.0 International licence, as recorded in the arXiv licence metadata for this exact version and shown on the abstract page https://arxiv.org/abs/2605.25146v1. Full rights notice: licenses/arxiv-2605.25146-CC-BY-4.0.txt.\par\smallskip

\noindent\textit{Editorial scope.} Counted unit: the Theorem whose source label is \texttt{thm:tens:lebesgue} in the article (source0.tex lines 1332--1343, source label \texttt{thm:tens:lebesgue}), delivered together with the article's own complete original proof of it (source0.tex lines 1345--1365, one \texttt{proof} environment of 21 lines). No mathematics is written, repaired or re-derived: statement and proof are the article's own text line for line. Required context retained uncounted: none. Every definition, display and label of those lines is retained unchanged. Omitted from this package and not redistributed: 1--1331, 1344--1344, 1366--2353. Nothing inside a retained excerpt is omitted or altered: every retained body is delivered exactly as the article's lines are written, so no omission receipt of any kind accompanies this package. Citations: the retained excerpts carry no citation command at all, so no bibliography entry of the article is transcribed and no reference is redistributed. Reference adaptation: none was needed and none was made; every reference inside the delivered unit resolves inside it, so no unresolved reference or citation is produced. Printed slips kept literal: no printed word, symbol, hypothesis or number of the counted statement or of its proof was altered. Layout and class: the article's own class is \texttt{imsart} with packages including geometry, fancyhdr, amsthm, amsmath, amsfonts, amssymb, dsfont, hyperref, cleveref, mathrsfs, graphicx, float, textcomp, rotating; the wrapper is the lane's pinned \texttt{amsart} editorial wrapper at 10pt on A4, which restates the theorem-like environments the retained text needs and adds the article's own macro definitions that the retained text needs (\texttt{\textbackslash b}, \texttt{\textbackslash c}, \texttt{\textbackslash innpr}, \texttt{\textbackslash m}, \texttt{\textbackslash mb}, \texttt{\textbackslash rd}, \texttt{\textbackslash s}, \texttt{\textbackslash vertiii}), with extra wrapper packages (bm, bbm, dsfont, xspace, cancel, mathtools). The delivered numbering is the unit's own. Assets: no retained span contains a figure environment, a graphics-inclusion command, a PDF-page inclusion, a table or a TikZ picture; no image file is copied into this package, the package declares no asset dependency and no image file is redistributed with it. Licence: version arXiv:2605.25146v1 of this article is distributed under the Creative Commons Attribution 4.0 International licence, as recorded in the arXiv licence metadata for this exact version and shown on the abstract page https://arxiv.org/abs/2605.25146v1. A full rights notice is delivered at licenses/arxiv-2605.25146-CC-BY-4.0.txt. Characters: every retained excerpt is pure ASCII, so the delivered engine, XeLaTeX with its default Latin Modern text font, reports no missing character.
\begin{theorem}[Lebesgue Dominated Convergence]\label{thm:tens:lebesgue}
Let $\c{Q}$ be a Hilbert space, and $\mb{\nu} \in \s{P}(X, \c{Q})$. An operator-valued measurable function $\m{F}: X \to (\c{B}_{\infty}(\c{H}), \mathrm{WOT})$ is tensorized Bochner integrable if there exists a sequence of measurable simple functions $\{\m{S}_{n}: X \to (\c{B}_{\infty}(\c{H}), \mathrm{WOT}) \}_{n \in \b{N}}$ such that:
\begin{enumerate}
    \item There exists a scalar function $h \in \c{L}_{1}(X, \mb{\nu})$ (meaning $h$ is integrable with respect to the scalar positive measures $\mb{\nu}_{\phi, \phi}$ for all $\phi \in \c{Q}$) satisfying $\vertiii{\m{S}_{n}(x)}_{\infty} \le h(x)$ for any $n \in \b{N}$ and $x \in X$.
    \item For each $x \in X$, $\m{S}_{n}(x) \to \m{F}(x)$ as $n \to \infty$ with respect to the WOT in $\c{B}_{\infty}(\c{H})$.
\end{enumerate}
In this case, we have $\vertiii{\m{F}(x)}_{\infty} \le h(x)$ for any $x \in X$, and
\begin{align*}
    \int_{X} \m{F}(x) \otimes \rd \mb{\nu}(x) = \lim_{n \to \infty} \int_{X} \m{S}_{n}(x) \otimes \rd \mb{\nu}(x),
\end{align*}
where the limit converges with respect to the WOT in $\c{B}_{\infty}(\c{H} \otimes \c{Q})$.
\end{theorem}

\begin{proof}
First, it is trivial that $\vertiii{\m{F}(x)}_{\infty} \le h(x)$ for any $x \in X$, since for any $f, g \in \c{H}$,
\begin{equation*}
    |\innpr{\m{F}(x) f}{g}_{\c{H}}| = \lim_{n \to \infty} |\innpr{\m{S}_{n}(x) f}{g}_{\c{H}}| \le h(x) \|f\|_{\c{H}} \|g\|_{\c{H}}.
\end{equation*}
Fix $f, g \in \c{H}$ and $\phi, \psi \in \c{Q}$. Because $|\innpr{\m{S}_{n}(x) f}{g}_{\c{H}}| \le h(x) \|f\|_{\c{H}} \|g\|_{\c{H}}$ and $h \in \c{L}_{1}(X, \mb{\nu}_{\phi, \psi})$, we can apply the standard Lebesgue dominated convergence theorem to obtain:
\begin{align*}
    \int_{X} \innpr{\m{F}(x) f}{g}_{\c{H}} \rd \mb{\nu}_{\phi, \psi} (x) &= \lim_{n \to \infty} \int_{X} \innpr{\m{S}_{n}(x) f}{g}_{\c{H}} \rd \mb{\nu}_{\phi, \psi} (x) \\
    &= \lim_{n \to \infty} \innpr{\left( \int_{X} \m{S}_{n}(x) \otimes \rd \mb{\nu}(x) \right) (f \otimes \phi)}{g \otimes \psi}_{\c{H} \otimes \c{Q}}.
\end{align*}
It suffices to show that the bounding limit is well-behaved for all $f, g \in \c{H}$ and $\phi, \psi \in \c{Q}$. Define the bounded positive operator $\m{H} := \int_{X} h(x) \rd \mb{\nu}(x) \in \c{B}_{\infty}^{+}(\c{Q})$, which is well-defined by the assumption $h \in \c{L}_{1}(X, \mb{\nu})$. By the Cauchy-Schwarz inequality for POVMs, the total variation of the complex measure satisfies $\rd |\mb{\nu}_{\phi, \psi}|(x) \le \sqrt{\rd\mb{\nu}_{\phi, \phi}(x) \rd\mb{\nu}_{\psi, \psi}(x)}$. 

Therefore, applying the Cauchy-Schwarz inequality for integrals, for any $n \in \b{N}$ we have:
\begin{align*}
    &\left| \int_{X} \innpr{\m{S}_{n}(x) f}{g}_{\c{H}} \rd \mb{\nu}_{\phi, \psi} (x) \right| 
    \le \|f\|_{\c{H}} \|g\|_{\c{H}} \int_{X} h(x) \rd |\mb{\nu}_{\phi, \psi}| (x) \\ 
    &\le \|f\|_{\c{H}} \|g\|_{\c{H}} \left( \int_{X} h(x) \rd \mb{\nu}_{\phi, \phi} (x) \right)^{1/2} \left( \int_{X} h(x) \rd \mb{\nu}_{\psi, \psi} (x) \right)^{1/2} \\
    &= \|f\|_{\c{H}} \|g\|_{\c{H}} \innpr{\m{H} \phi}{\phi}_{\c{Q}}^{1/2} \innpr{\m{H} \psi}{\psi}_{\c{Q}}^{1/2} \le \|f\|_{\c{H}} \|g\|_{\c{H}} \vertiii{\m{H}}_{\infty} \|\phi\|_{\c{Q}} \|\psi\|_{\c{Q}}.
\end{align*}
This uniform bound over all $n$ strictly guarantees that the sesquilinear form defines a bounded operator on the tensor product space, completing the proof.
\end{proof}

\end{document}
